\documentclass[11pt]{article}

% A clean starting point for a new manuscript. Everything here compiles with
% OffLeaf's defaults (pdfLaTeX via latexmk). Replace the placeholder text and
% go. Delete any packages you don't need.
\usepackage[utf8]{inputenc}
\usepackage{amsmath, amssymb}   % mathematics
\usepackage{graphicx}           % \includegraphics
\usepackage[numbers]{natbib}    % citations: \citep{key}, \citet{key}
\usepackage{hyperref}           % clickable links and references

\title{Your Title Here}
\author{Your Name}
\date{\today}

\begin{document}
\maketitle

\begin{abstract}
One paragraph summarizing the problem, the approach, and the headline result.
Aim for 150--250 words.
\end{abstract}

\section{Introduction}
\label{sec:intro}
State the problem and why it matters. Cite prior work like
this~\citep{example2024}. Inline math works anywhere: $f(x) = x^2$.

\section{Methods}
\label{sec:methods}
Describe what you did, precisely enough to reproduce. Display equations get
numbers you can reference:
\begin{equation}
  \Delta G = -k_\mathrm{B} T \ln K,
  \label{eq:deltaG}
\end{equation}
and you can refer back to Eq.~\eqref{eq:deltaG} or Section~\ref{sec:intro}.

\section{Results}
\label{sec:results}
What you found. Figures go in a \texttt{figs/} folder next to this file:
% \begin{figure}[htbp]
%   \centering
%   \includegraphics[width=0.8\linewidth]{figs/myplot.png}
%   \caption{What the figure shows.}
%   \label{fig:myplot}
% \end{figure}

\section{Discussion}
\label{sec:discussion}
What it means, limitations, and what comes next.

\bibliographystyle{unsrtnat}
\bibliography{refs}

\end{document}
